\documentclass[11pt]{article}
\usepackage[margin=1in]{geometry}
\usepackage{amsmath,amssymb,amsthm}
\usepackage{graphicx}
\usepackage{grffile}
\usepackage[onehalfspacing]{setspace}
\usepackage[unicode=true,
 bookmarks=true,bookmarksnumbered=false,bookmarksopen=false,
 breaklinks=false,pdfborder={0 0 0},pdfborderstyle={},backref=false,colorlinks=true]
 {hyperref}
\usepackage{verbatim}
\usepackage{appendix}
\usepackage{mathtools}
\usepackage{url}
\usepackage[authoryear]{natbib}
\usepackage{booktabs} %Package with \toprule, \midrule, \bottomrule
\usepackage{placeins} %Package with \FloatBarrier command
%\usepackage{physics}
\usepackage{datetime}
\usepackage{epstopdf}
\usepackage{changepage}
\usepackage{subfig}
\usepackage[normalem]{ulem} % Customize underline
\captionsetup[subfigure]{labelformat=simple, labelsep=period}
\renewcommand{\thesubfigure}{\Alph{subfigure}}
\usepackage{etoolbox}  %Supports toggle commands
\newcommand{\appendixpagenumbering}{ 
	\pagenumbering{arabic}
	\renewcommand{\thepage}{Appendix - \arabic{page}}
}
\newtoggle{12ptfont}
\togglefalse{12ptfont}

\hypersetup{
	pdftex,
    pdftitle={Spatial Economics for Granular Settings},
    pdfauthor={Jonathan I. Dingel and Felix Tintelnot},
    colorlinks=true,       % false: boxed links; true: colored links
    linkcolor=red,          % color of internal links
    citecolor=black,        % color of links to bibliography
    urlcolor=blue           % color of external links
}

\hyphenpenalty=500

% CUSTOM DEFINITIONS
\urldef{\dingelemail}\url{j.dingel@columbia.edu}
\urldef{\tintelnotemail}\url{felix.tintelnot@duke.edu}
\newtheorem{proposition}{Proposition}
\theoremstyle{definition}
\newtheorem{definition}{Definition}[section]
\newtheorem*{remark}{Remark}

% IDENTIFYING INFORMATION
\title{Spatial Economics for Granular Settings\thanks{
	We thank four anonymous referees,
	Rodrigo Ad\~{a}o, Gabriel Ahlfeldt, Treb Allen, Stephane Bonhomme, Kirill Borusyak, Victor Couture, Teresa Fort, Lars Hansen, Elhanan Helpman, Thomas Holmes, John Huizinga, Erik Hurst, Kyle Mangum, Yuhei Miyauchi, Eduardo Morales, Stephen Redding, Jordan Rosenthal-Kay, Esteban Rossi-Hansberg, Chris Severen, Sebastian Sotelo, Daniel Sturm, Miguel Zerecero, Oren Ziv, and
	numerous conference and seminar participants for helpful feedback.
	We are grateful to Junbiao Chen, Levi Crews, Gianna Fenaroli, Reigner Kane, Xianglong Kong, Daniil Iurchenko, Leran Qi, John Ruf, Ye Sun, Linghui Wu, Shijian Yang, and Mingjie Zhu for excellent research assistance.
	Dingel thanks the James S. Kemper Foundation Faculty Research Fund at the University of Chicago Booth School of Business.
	This work was completed in part with resources provided by the University of Chicago Research Computing Center
	and Columbia University's Shared Research Computing Facility.
	We gratefully acknowledge support from the National Science Foundation under grant number 2018609.
	\dingelemail \ and \tintelnotemail
}}
\author{Jonathan I. Dingel
\\ Columbia University, NBER, and CEPR
\and
Felix Tintelnot
\\ Duke University, NBER, and CEPR}
\date{\the\day \ \monthname[\the\month] \the\year}

\newcommand\idiosyncraticphrase{idiosyncratic preference}

\begin{document}
\begin{figure}[hbt]
\caption{Number of commuters between pairs of tracts in New York City}
\label{fig:NYC:histogram}
\vspace{-5mm}
\begin{center}
\includegraphics[width=0.60\textwidth]{../input/histogram_NYC_2010_positive_w_shares.eps}
\begin{minipage}{1.0\textwidth}
{\footnotesize \textsc{Notes}: 
This figure describes tract pairs in New York City by the number of individuals who reside in the origin tract and work in the destination tract in 2010 LODES data.
The sample is restricted to the 15\% of tract pairs that have a strictly positive number of commuters.
The blue histogram depicts the fraction of tract pairs reporting each number of commuters,
for each value from 1 to 39 and values of 40 or greater.
The black increasing step function depicts the cumulative share of commuters by the number of commuters between the pair of tracts.
New York City has 2,160 residential census tracts.%
\par}
\end{minipage}
\end{center}
\end{figure}
\begin{figure}[hbt]
\caption{Calibrated-shares procedure overfits in Monte Carlo simulations}
\label{fig:montecarlo:MSEanddensities}
\begin{center} \vspace{-4mm}
\subfloat[Regression of true on predicted changes]{
	\begin{minipage}{0.4\textwidth} Exhibit omitted from ``fast'' pipeline. Monte Carlo simulations are slow. \end{minipage}
}
\subfloat[Ratio of models' prediction errors]{
	\begin{minipage}{0.4\textwidth} Exhibit omitted from ``fast'' pipeline. Monte Carlo simulations are slow. \end{minipage}
}
\\
\begin{minipage}{0.95\textwidth}
{\footnotesize \textsc{Notes}:
This figure depicts the regression coefficients and mean squared errors from 100 simulations in which $I = 2,488,905$,
the data-generating process is the estimated covariates-based
model of New York City in 2010,
and the counterfactual change is a 9\% increase in productivity in one workplace tract.
See Appendix~A.1 for a detailed description of the simulation procedure.
``Covariates-based'' means estimating the covariates-based parameterization described in Section~3.2 using the simulated data.
``Calibrated-shares'' means using the observed (simulated) shares in the exact hat algebra defined by equations~(5)--(7).
In Panel~A, the densities for the covariates-based model are on the left vertical axis and the densities for the calibrated-shares procedure are on the right vertical axis.
\par}
\end{minipage}
\end{center}
\end{figure}
\begin{figure}[!h]
\caption{Comparison of models' predictive performance across 83 events}
\label{fig:event:83events:coefficientdensities}
\begin{center}\vspace{-4mm}
\subfloat[Regression of observed on predicted changes]{ 
\includegraphics[width=.49\linewidth]{../input/cbm_sigma_4.0_vs_csp_sigma_4.0_kdensity_exhibit.eps}
\label{fig:eventstudy_coef_cont_cs}
}
\subfloat[Ratio of models' prediction errors]{
\includegraphics[width=.49\linewidth]{../input/cbm_sigma_4.0_vs_csp_sigma_4.0_histogram_exhibit.eps}\label{fig:mse_cont_cs}
}\\ \vspace{-2mm}
\subfloat[No extensive margin: Regression coefficients]{
\includegraphics[width=.49\linewidth]{../input/cbm_deltainf_vs_csp_sigma_4.0_kdensity_exhibit.eps}
\label{fig:eventstudy_coef_continfzero_cs}}
\subfloat[No extensive margin: Ratio of prediction errors]{
\includegraphics[width=.49\linewidth]{../input/cbm_deltainf_vs_csp_sigma_4.0_histogram_exhibit.eps}\label{fig:mse_cc}
} \\
\begin{minipage}{0.98\linewidth}{\footnotesize \textsc{Notes}: 
``Calibrated-shares'' refers to using the observed shares in the exact hat algebra defined by equations~(5)--(7).
``Covariates-based'' in Panels~\subref*{fig:eventstudy_coef_cont_cs} and \subref*{fig:mse_cont_cs} refers to the parameterization described in Section~3.2.
``No extensive margin'' in Panels~\subref*{fig:eventstudy_coef_continfzero_cs} and \subref*{fig:mse_cc}
retains only the intensive-margin predictions of the covariates-based model.
This case uses 
the covariates-based model's predictions for residence-workplace pairs with $\ell_{kn}>0$ in the 2010 data
and 
predicts zero change for pairs with $\ell_{kn}=0$ in the 2010 data.
In Panel~\subref*{fig:eventstudy_coef_cont_cs},
\input{../input/cbm_sigma_4.0_vs_csp_sigma_4.0_kdensity_notes.tex}
In Panel~\subref*{fig:eventstudy_coef_continfzero_cs},
3 slope and 1 intercept coefficients are not depicted for the calibrated-shares procedure.
Vertical lines depict the medians of the slope coefficient distributions.
\par}
\end{minipage}
\end{center}
\end{figure}
\begin{figure}[!tbh]
\caption{Predictive performance across 83 events using low-rank approximation}
\label{fig:event:83events:SVD:coefficientdensities}
\begin{center}\vspace{-4mm}
\subfloat[Regression of observed on predicted changes]{
\label{fig:density_slope_SVD}
\includegraphics[width=.45\textwidth]{../input/cbm_sigma_4.0_vs_svd_16_kdensity_exhibit.eps}
}
\subfloat[Ratio of models' prediction errors]{
\label{fig:hist_MSE_SVD}
\includegraphics[width=.45\textwidth]{../input/cbm_sigma_4.0_vs_svd_16_histogram_exhibit.eps}
}\\ \vspace{-2mm}
\begin{minipage}{0.90\textwidth}{\footnotesize
\textsc{Notes}:
This figure depicts the predictive performance of 
using an approximation of the commuting matrix in exact hat algebra.
The approximation is derived from a rank-16 singular value decomposition of the observed commuting matrix,
as described in Appendix~B.4.
The covariates-based model is the same as in Figure~\ref{fig:event:83events:coefficientdensities}.
Panel~\subref*{fig:density_slope_SVD} depicts the distributions of slope and intercept coefficients, and the vertical lines denote the median slope coefficients.
\par}\end{minipage}\\ \vspace{4mm}
\end{center}
\end{figure}
\begin{figure}[!h]
\caption{Predictive performance of the rank-1 IFE specification across 83 events}
\label{fig:event:83events:IFE:coefficientdensities}
\begin{center}\vspace{-4mm}
\subfloat[Regression of observed on predicted changes]{
\label{fig:density_slope_IFE}
\begin{minipage}{0.40\textwidth} Exhibit omitted from ``fast'' pipeline. IFE estimation is slow due to many starting point in numerical optimization.\end{minipage}
}
\subfloat[Ratio of models' prediction errors]{
\label{fig:hist_MSE_IFE}
\begin{minipage}{0.40\textwidth} Exhibit omitted from ``fast'' pipeline. IFE estimation is slow due to many starting point in numerical optimization.\end{minipage}
}\\ \vspace{-2mm}
\begin{minipage}{0.90\textwidth}{\footnotesize \textsc{Notes}: 
This figure summarizes the predictive performance of the covariates-based model 
and the rank-1 IFE specification.
Panel~\subref*{fig:density_slope_IFE} depicts the distributions of slope coefficients, and the vertical lines denote the median slope coefficients.
\par}\end{minipage}\\ \vspace{4mm}
\end{center}
\end{figure}
\begin{table}[hbt]
\caption{Commuting elasticity estimates}
\label{tab:elasticity}
\begin{center}
\vspace{-5mm}
\begin{tabular}{lcc} \toprule
&PPML/MLE & OLS \\ \midrule
Commuting cost & -7.986 & -2.307 \\
 & (0.307) & (0.00768) \\
 &  &  \\
Model fit ($R^2$ or pseudo-$R^2$) & 0.662 & 0.561 \\
Location pairs & 4,628,878 & 690,673 \\
 Commuters & 2,488,905 & 2,488,905 \\ \bottomrule
\end{tabular}

\begin{minipage}{1.0\textwidth}
{\footnotesize \textsc{Notes}: 
All specifications include residence fixed effects and workplace fixed effects.
The ``PPML/MLE'' column presents the results from maximum likelihood estimation of equation~(8).
The ``OLS'' column presents the results of estimating the log version of equation~(2) by ordinary least squares, omitting observations in which $\ell_{kn}=0$.
The model-fit statistic is the pseudo-$R^2$ for MLE and $R^2$ for OLS.
We report the PPML standard errors (clustered by $k$ and by $n$), which are larger than the logit MLE standard errors associated with maximizing equation~(8).
\par}
\end{minipage}
\end{center}
\end{table}
\begin{table}[tbh]
\caption{Calibrated-shares procedure's finite-sample performance}
\label{table:montecarlo:cs_lambdazero:maintext}
\vspace{-4mm}
\begin{center}
A. Regressand is continuum change in commuters\\
\resizebox{0.6\textwidth}{!}{%
Exhibit omitted from ``fast'' pipeline. Monte Carlo simulations are slow.
}\\
B. Regressand is finite-sample change in commuters\\
\resizebox{0.6\textwidth}{!}{%
Exhibit omitted from ``fast'' pipeline. Monte Carlo simulations are slow.
}
\iftoggle{12ptfont}{\begin{minipage}{1.0\textwidth}}{\begin{minipage}{0.95\textwidth}}
{\footnotesize \textsc{Notes}:
This table reports the average values of the slope coefficient, intercept coefficient, and mean squared error from 100 simulations as we vary $I$, the number of individuals in the simulated economy.
The column titles denote the number of individuals in millions.
The ``2.5'' million case corresponds to $I = 2,488,905$,
which is the number of individuals who reside and work in New York City in the 2010 LODES data.
In Panel~A, the regressand is the change in commuters with a continuum of individuals.
This change does not vary across simulations.
In Panel~B, the regressand is the simulation-specific change in commuters from a realization drawn from the model using pre-shock parameter values to a realization drawn from the model using post-shock parameter values.
\par}
\end{minipage}
\end{center}
\end{table}
\begin{figure}[hbt]
\caption{Amazon HQ2 counterfactual change in residents}
\begin{center}\vspace{-4mm}
\subfloat[Residents working in HQ2 tract]{ 
\includegraphics[width=.33\linewidth]{../input/map_res2010.png}
\label{fig:NYC:AmazonHQ2:origins2010}}
\subfloat[Covariates-based predictions]{ 
\includegraphics[width=.33\linewidth]{../input/map_cont_reschange.png}\label{fig:NYC:AmazonHQ2:CBM:residentchange}
}
\subfloat[Calibrated-shares predictions]{
\includegraphics[width=.33\linewidth]{../input/map_cs_changeres.png}\label{fig:NYC:AmazonHQ2:CSP:residentchange}
}\\
\begin{minipage}{\textwidth}{\footnotesize
\textsc{Notes}:
Panel~\subref*{fig:NYC:AmazonHQ2:origins2010} depicts the number of residents
in each tract who work in the treated Long Island City tract in the 2010 LODES data.
The legend of Panel~\subref*{fig:NYC:AmazonHQ2:origins2010} reports
the percentiles corresponding to the integer number of residents,
excluding the treatment tract.
Panels~\subref*{fig:NYC:AmazonHQ2:CBM:residentchange} and \subref*{fig:NYC:AmazonHQ2:CSP:residentchange} depict changes in the number of residents predicted by the covariates-based model and calibrated-shares procedure, respectively.
The legend percentile cutoffs in Panels~\subref*{fig:NYC:AmazonHQ2:CBM:residentchange} and \subref*{fig:NYC:AmazonHQ2:CSP:residentchange} correspond to those in Panel~\subref*{fig:NYC:AmazonHQ2:origins2010}.
The covariates-based predictions for residents describe both the continuum model and the model with finite individuals
because the change in the former equals the expected change in the latter.
\par}
\end{minipage}
\end{center}
\end{figure}
\begin{figure}[bth] 
\caption{Predicted changes in rents}
\label{fig:AHQ2:rentchange}
\begin{center}\vspace{-5mm}
\subfloat[Covariates-based: Continuum]{
\includegraphics[width=.33\linewidth]{../input/map_cont_change_realr.png}
\label{fig:AHQ2:rentchange:CBM}}
\subfloat[Covariates-based: Finite]{
\begin{minipage}{0.3\textwidth} Exhibit omitted from ``fast'' pipeline. Fixed $\nu$ simulations are slow. \end{minipage}
\label{fig:AHQ2:rentchange:CBM:finitemodel}}
\subfloat[Calibrated shares]{
\includegraphics[width=.33\linewidth]{../input/map_cs_hat_realr.png}
\label{fig:AHQ2:rentchange:CSP}}\\
\begin{minipage}{0.90\textwidth}{\footnotesize
\textsc{Notes}:
These maps depict percent changes in rents, $(\hat{r}_{k}/\hat{P} - 1) \times 100$.
Panel~\subref*{fig:AHQ2:rentchange:CBM}
depicts the change in the covariates-based continuum model.
Panel~\subref*{fig:AHQ2:rentchange:CBM:finitemodel}
depicts the mean change across 100 simulations of the model with finite individuals.
Panel~\subref*{fig:AHQ2:rentchange:CSP}
depicts each tract's predicted rent change using the calibrated-shares procedure.
\par}
\end{minipage}
\end{center}
\end{figure}
\begin{figure}[h!]
\caption{Variation in counterfactual changes: Idiosyncrasies vs. parameter uncertainty}
\label{fig:granularity:fixednu}
\label{fig:NYC:AmazonHQ2:parameteruncertainty}
\begin{center}\vspace{-4mm}
\subfloat[Changes in residents]{ 
	\begin{minipage}{0.3\textwidth} Exhibit omitted from ``fast'' pipeline. Fixed $\nu$ simulations are slow. \end{minipage}
\label{fig:AHQ2:puncertainty:res}\label{fig:fixednu:g_res_simulation}}
\subfloat[Changes in workers]{
	\begin{minipage}{0.3\textwidth} Exhibit omitted from ``fast'' pipeline. Fixed $\nu$ simulations are slow. \end{minipage}
\label{fig:AHQ2:puncertainty:emp}\label{fig:fixednu:g_emp_simulation}}\\ \vspace{-2mm}
\subfloat[Changes in rents]{
	\begin{minipage}{0.3\textwidth} Exhibit omitted from ``fast'' pipeline. Fixed $\nu$ simulations are slow. \end{minipage}
\label{fig:AHQ2:puncertainty:realrent}\label{fig:fixednu:g_realrent}}
\subfloat[Changes in wages]{
	\begin{minipage}{0.3\textwidth} Exhibit omitted from ``fast'' pipeline. Fixed $\nu$ simulations are slow. \end{minipage}
\label{fig:AHQ2:puncertainty:realwage}\label{fig:fixednu:g_realwage}}\\ \vspace{-2mm}
\begin{minipage}{0.96\textwidth}{\footnotesize
\textsc{Notes}:
These plots depict variation in counterfactual changes 
induced by individual idiosyncrasies in the model with a finite number of individuals
alongside
variation attributable to parameter uncertainty about $\{A_n\}$, $\{T_k\}$, and $\epsilon$ in the covariates-based continuum model.
We summarize the former by plotting the 5\textsuperscript{th} and 95\textsuperscript{th} percentiles of predicted changes in quantities and prices across 100 simulations of the model with a finite number of individuals.
In Panels~\subref*{fig:fixednu:g_res_simulation}, \subref*{fig:fixednu:g_emp_simulation}, \subref*{fig:fixednu:g_realrent}, and \subref*{fig:fixednu:g_realwage},
the horizontal axes correspond to the values depicted in 
Figures~\ref{fig:NYC:AmazonHQ2:CBM:residentchange}, F.1B, \ref{fig:AHQ2:rentchange:CBM:finitemodel}, and F.2B, respectively.
Panels \subref*{fig:AHQ2:puncertainty:emp} and \subref*{fig:AHQ2:puncertainty:realwage} exclude the Amazon HQ2 workplace tract.
\dots
We quantify parameter uncertainty by producing 100 counterfactual predictions by
sampling with replacement 2,488,905 individuals from the observed commuting matrix 100 times.
For each of these 100 samples,
we estimate the covariates-based continuum model and compute its counterfactual predictions.
The plots depict the 5\textsuperscript{th} and 95\textsuperscript{th} percentiles of predicted changes in quantities and prices across these 100 bootstrap samples.
The variation attributable to parameter uncertainty is considerably smaller 
than 
the variation attributable to individual idiosyncrasies:
the bootstrap 5\textsuperscript{th} and 95\textsuperscript{th} percentiles are almost indistinguishable.
\par}
\end{minipage}
\end{center}
\end{figure}
\end{document}
